% Guideline item 4: Motivation.
\section{Motivation}\label{sec:motivation}

\paragraph{A worked example: the one-way ratchet.}
Take a threshold that keeps the last $L=3$ admitted scores, sets itself to the largest
of them, and admits a new score \emph{only if it judged that score normal}
($s \le \theta$). It starts at $\theta=10$ with an empty buffer.
Table~\ref{tab:ratchet-example} follows six scores.

\begin{table}[H]
  \centering
  \caption{Six steps of a self-referential threshold ($L=3$, threshold = buffer
  maximum). A score of 9 is normal at step 1 and an alert at step 6.}
  \label{tab:ratchet-example}
  \small
  \begin{tabular}{@{}cccccc@{}}
    \toprule
    Step & Score $s$ & Judged against & Decision & Buffer after & Threshold after\\
    \midrule
    1 & 9  & 10 & normal, admitted  & \{9\}      & 9\\
    2 & 12 & 9  & alert, rejected   & \{9\}      & 9\\
    3 & 8  & 9  & normal, admitted  & \{9, 8\}   & 9\\
    4 & 7  & 9  & normal, admitted  & \{9, 8, 7\} & 9\\
    5 & 6  & 9  & normal, admitted  & \{8, 7, 6\} & 8\\
    6 & 9  & 8  & \textbf{alert}, rejected & \{8, 7, 6\} & 8\\
    \bottomrule
  \end{tabular}
\end{table}

Every admitted score is at most the current threshold, so the buffer maximum can only
stay the same or fall. Over a long run the threshold sinks toward the bottom of the
normal band, and ordinary fluctuations above it become alerts. The alert flood says
nothing about the signal; it is produced by a threshold that learns only from what it
already called normal. Section~\ref{sec:problem} proves this for any score sequence,
buffer length and quantile (Proposition~\ref{prop:ratchet}).

\paragraph{The opposite reflex fails too.}
Admitting \emph{every} score removes the ratchet but creates the reverse failure. When a
sustained fault arrives, its scores fill the buffer, the threshold rises to meet them,
and the fault becomes the new normal. One reflex cannot rise; the other cannot tell a
fault from a new normal. Our question sits between them.

\paragraph{Why this is not already solved.}
Drift-adaptive detectors retrain the \emph{model} when drift is
detected~\citep{yang2026ddade, li2026adapts, xu2024cddia}, abstention methods withhold
decisions statistically~\citep{sun2024online, perini2023rejection}, and FITNESS already
reports that discarding self-declared anomalies can fail~\citep{sankararaman2022fitness}.
None of the nine closest papers we checked treats the threshold's choice of training
data as an explicit action and compares that action with its alternatives on one fixed
score trace, with alert workload and decision coverage reported apart. Without that
common trace a reported gain can come from a better scorer rather than a better
decision, which is the attribution problem this study is built to avoid.
